\documentclass{article}

\usepackage[utf8]{inputenc}
\usepackage[T1]{fontenc}

\usepackage{amsmath}
\usepackage{amssymb}
\usepackage{amsfonts}
\usepackage{textcomp}
\usepackage{pifont}
\usepackage{gensymb}

\usepackage{listings}

\title{Resumo}

\begin{document}

simbologia -> começa depois de sobrescritos e subscritos

símbolos de operações simples
\begin{lstlisting}
a + b
a - b
a \cdot b
\end{lstlisting}
\[
  a + b
  \quad
  a - b
  \quad
  a \cdot b
\]

símbolos de operações simples 2
\begin{lstlisting}
\frac{}{}
\sqrt{}
(...) [...] \left ( ... \right )
\end{lstlisting}
\[
  \frac{x+1}{x-1} \qquad
  \sqrt[n]{2} \qquad
  (x+1) \quad [x+1] \quad\left ( \frac{x+1}{x-1} \right)
\]

símbolos gregos

equações longas -> depois de porque usar $\sin$

matrizes
\begin{lstlisting}
\begin{matrix}
a & b
\\
c & d
\end{matrix}
pmatrix para parenteses
bmatrix para colchetes
vetores
indices com _{}
\end{lstlisting}
\[
  \begin{matrix}
  a & b
  \\
  c & d
  \end{matrix}
  \quad
  \begin{pmatrix}
  a & b
  \\
  c & d
  \end{pmatrix}
  \quad
  \begin{bmatrix}
  a & b
  \\
  c & d
  \end{bmatrix}
  \quad
  \begin{pmatrix}
  a
  \\
  b
  \end{pmatrix}
  \quad
  A =
  \begin{pmatrix}
  a_{11} & a_{12}
  \\
  a_{21} & a_{22}
  \end{pmatrix}
\]

equações
\begin{lstlisting}
& define ponto de alinhamento
\\ quebra a linha
align
cases
\end{lstlisting}

  \begin{align*}
  x + y &= 10
  \\
  x-y &= 2
  \end{align*}
\[
  \quad
  \begin{cases}
  x + y = 10
  \\
  x - y = 2
  \end{cases}
  \quad
  f(x) = 
  \begin{cases}
  x^2, & x\geq 0
  \\
  -x, & x<0
  \end{cases}
\]



\end{document}